\section{Conclusion}

The question posed at the outset --- what would it take, experimentally, to establish that an ecological system carries fractional memory --- receives a quantitative and partly negative answer. Four analytical results in the main text (structural separation, finite-horizon latent approximation, the complexity-dependent testing obstruction, and the interval-certified prey-response approximation of Theorem~\ref{thm:T23}) delimit what is provable; the numerical layers (frequency diagnostics, the fractional-delay study of Section~\ref{sec:fd}, six input families with safety metrics, and the 1512-cell benchmark) measure what is achievable.

The measured answer is consistent across layers. Fractional memory is structurally distinct, yet at latent budget $m=32$ the certified surrogate is within $0.002$ of chance-level indistinguishability under the declared pulse protocol; excitation strong enough to separate the mechanisms (macro-accuracy up to $0.584$) crosses the Allee threshold in up to $100\%$ of stable cells. Before the next experiment, the productive question is therefore not \emph{which model fits best} but \emph{which safe perturbation makes the competing memory mechanisms observably different}.
